El diseño metodológico se estructura en cuatro fases secuenciales, cada una con criterios de validez, instrumentos y entregables definidos.

\section{Aspectos principales de la metodología}

El proyecto adopta un enfoque experimental cuantitativo con diseño cuasi-experimental, donde la variable independiente es la configuración del sistema (LLM solo vs. sistema híbrido LLM + agente) y las variables dependientes son las métricas de desempeño (F1, precisión, recall) y explicabilidad. Se utilizan múltiples modelos de lenguaje y configuraciones de prompting para controlar el efecto del modelo específico sobre los resultados.

\textbf{Hipótesis:} El sistema híbrido que combina un LLM clasificador con un agente verificador de fuentes institucionales colombianas alcanza un F1 macro significativamente superior al LLM utilizado de forma independiente en la detección de desinformación financiera colombiana.

\section{Método de investigación}

\subsection{Datos}

\subsubsection{Universo y muestra}

El universo de estudio comprende las noticias financieras y bancarias que circulan en medios digitales colombianos y que han sido objeto de verificación por organizaciones certificadas por la IFCN, así como noticias publicadas por medios económicos de referencia. La muestra objetivo es de 800 a 1.000 noticias, distribuidas en tres clases.

\textbf{Nota sobre fuentes (actualizado durante la ejecución de Fase 1):} la lista de fuentes descrita originalmente en el anteproyecto incluía Portafolio, El Tiempo y AFP Factual Colombia. Al verificar técnicamente cada fuente (ver \texttt{03\_scraping/README.md}), se encontró que Portafolio y El Tiempo bloquean explícitamente bots de IA en su \texttt{robots.txt} (nombran \texttt{anthropic-ai} y \texttt{ClaudeBot} por nombre) y que AFP Factual bloquea activamente cualquier solicitud automatizada a nivel de firewall de aplicación (Akamai). Las tres quedaron excluidas -- no por decisión metodológica, sino porque scrapearlas violaría la propia regla de este documento de respetar \texttt{robots.txt}. Se compensó agregando dos medios de referencia adicionales verificados como técnicamente accesibles (Halcones y Palomas, Valora Analitik), y las clases finales usadas son:

\begin{itemize}
    \item \textbf{Verdaderas (300--400):} La República, Halcones y Palomas, Valora Analitik (medios de referencia), y comunicados oficiales del Banco de la República y la Superintendencia Financiera de Colombia.
    \item \textbf{Falsas (300--400):} Noticias calificadas como ``Falso'' por ColombiaCheck, único fact-checker certificado IFCN que resultó técnicamente accesible.
    \item \textbf{Dudosas (200--300):} Noticias calificadas por ColombiaCheck como ``Cuestionable'' o ``Verdadero, pero'' (taxonomía real verificada en \texttt{colombiacheck.com/metodologia}, ver Estado del Arte), así como contenido viral financiero sin verificación institucional.
\end{itemize}

El DANE quedó pendiente: sus comunicados de prensa con cifras económicas (IPC, PIB) son mayoritariamente PDF, no HTML, lo que requiere un extractor distinto al usado para el resto de fuentes (Trafilatura, que opera sobre HTML) -- no implementado en esta fase.

\subsubsection{Adquisición}

La recolección se implementa como un proyecto Scrapy (\texttt{03\_scraping/}), con Trafilatura para la extracción de texto limpio a partir del HTML crudo de cada artículo -- se prefirió Trafilatura sobre selectores CSS por sitio porque detecta el contenido principal de forma genérica, sin depender de la estructura particular de cada página. El pipeline de procesamiento aplica, en orden, sobre cada ítem candidato: (1) un filtro de dominio financiero basado en palabras clave (recall amplio, deliberadamente simple para no introducir un clasificador LLM en la etapa de recolección y contaminar la evaluación de Fase 2); (2) una lista de exclusión manual persistida en código (\texttt{scraper/exclusiones\_manuales.py}), resultado de una calibración inicial donde se revisaron a mano 58 ítems de ColombiaCheck con coincidencia incidental de una sola palabra clave -- 53 resultaron ser ruido temático (contenido político con una mención financiera incidental) y se excluyeron permanentemente; (3) deduplicación por similitud de n-gramas de caracteres (umbral $>0.85$), con el índice sembrado con \emph{todo} lo ya recolectado de las demás fuentes, no solo de la corrida en curso, para detectar la misma noticia reportada por dos medios distintos; y (4) validación de esquema, incluyendo que un ítem sin fecha de publicación verificable se descarta, no se etiqueta con una fecha inventada.

Cada ítem recibe un identificador determinístico: \texttt{id = SHA-256(url)[:16]}. Esto tiene dos propiedades relevantes para la reproducibilidad de la limpieza de datos: la misma URL produce siempre el mismo id entre corridas distintas (permite deduplicar y hacer seguimiento de un ítem específico sin depender de un contador incremental que cambiaría de corrida a corrida), y el id no requiere una base de datos central ni coordinación entre spiders para evitar colisiones.

El resultado de cada spider se escribe en formato \textbf{Parquet} (\texttt{02\_corpus/raw/\{fuente\}.parquet}), un formato columnar comprimido, elegido sobre CSV por preservar tipos de datos (fechas, booleanos) sin ambigüedad de parseo y por su eficiencia de lectura con pandas/PyArrow para los volúmenes del proyecto. El esquema completo de columnas está documentado en \texttt{02\_corpus/README.md}. Los datos crudos no se versionan en git (son datos pesados y algunos artículos tienen derechos de autor de terceros, ver Marco Normativo); el código que los produce sí. Cada corrida escribe además una copia en \textbf{CSV} (con BOM UTF-8) en paralelo al Parquet, exclusivamente para inspección humana directa (abrir en Excel, revisión de avance con el director) -- Parquet sigue siendo el formato que consume el resto del pipeline. El avance de recolección se resume en \texttt{06\_evaluacion/avance\_corpus.md}, regenerable con \texttt{06\_evaluacion/generar\_avance\_corpus.py}.

Se respetan estrictamente los archivos \texttt{robots.txt} de cada sitio (\texttt{ROBOTSTXT\_OBEY=True} a nivel de configuración del proyecto, sin excepciones) y se usa un \textit{User-Agent} que identifica honestamente el propósito académico del scraper, en lugar de simular un navegador para evadir bloqueos.

Como mecanismo complementario, se integrarán datasets públicos en español (FakeDeS/IberLEF 2021, 971 ítems, re-etiquetados según la taxonomía trinaria cuando pertenezcan al dominio financiero) y se dejó diseñado -- no implementado todavía -- un pipeline periódico vía GitHub Actions para recolección incremental.

\subsubsection{Reproducibilidad de la recolección}

Es importante declarar explícitamente qué parte de este proceso es reproducible y cuál no. El \textbf{código} es completamente reproducible: está versionado en git, es determinístico dado el mismo HTML de entrada (mismo filtro, mismo umbral de deduplicación, mismas reglas de exclusión ya persistidas), y cualquier persona con acceso al repositorio puede ejecutar los mismos spiders con la misma configuración.

El \textbf{corpus resultante no es perfectamente reproducible}, y esto es una limitación inherente a cualquier corpus construido por scraping de sitios web en vivo, no un defecto específico de este proyecto: los sitios fuente publican contenido nuevo y en ocasiones retiran o editan contenido antiguo, de modo que ejecutar el mismo scraper en una fecha distinta no produce exactamente el mismo conjunto de noticias. Se mitiga este problema de dos formas: (1) cada ítem registra su \texttt{fecha\_recoleccion} además de su \texttt{fecha\_publicacion}, de modo que cualquier resultado reportado en esta tesis puede vincularse a una versión fechada específica del corpus, no a ``lo que el scraper produzca en cualquier momento''; y (2) el corpus final usado para los resultados reportados (no solo el código para generarlo) se conservará como un artefacto versionado y fechado, para que la evaluación de Fase 2 y 3 sea reproducible contra esa versión específica, aunque el sitio origen cambie después.

\textbf{Tratamiento y trazabilidad de correcciones de extracción.} Durante la recolección se identificaron errores puntuales de extracción (texto contaminado con elementos decorativos de la página, calificaciones de fact-checker no contempladas en el mapeo inicial). Cada corrección de este tipo se documenta en \texttt{03\_scraping/REGISTRO\_CORRECCIONES.md} con fecha, causa raíz, número de ítems afectados y el cambio de código aplicado -- ninguna corrección edita un dato ya extraído a mano; todas corrigen el código de extracción y regeneran los datos desde el HTML ya almacenado en caché. Esto es deliberado: sin este registro, una diferencia en los conteos del corpus entre corridas sucesivas no tendría forma de distinguirse de una manipulación de datos no declarada.

\subsubsection{Muestreo final del corpus}

La recolección descrita arriba produce un volumen mayor al objetivo (más de 1.200 ítems frente a 800-1.000), y con una composición de fuentes dispar dentro de cada clase -- por ejemplo, dentro de la clase ``verdadera'', una sola fuente puede llegar a representar más del 40\% de los ítems, simplemente porque esa fuente tenía más contenido disponible para scrapear, no porque sea más representativa del universo real de noticias financieras colombianas. Reducir el corpus recolectado al tamaño objetivo mediante \textbf{muestreo aleatorio simple o estratificado proporcional} preservaría ese desbalance intacto (un muestreo proporcional, por definición, mantiene las proporciones de la población de origen).

Por eso el muestreo final usa \textbf{muestreo estratificado con asignación desproporcionada}~\cite{Cochran1977}: los estratos se definen por la combinación clase × fuente (no solo por clase), y la asignación de cuota por estrato no es proporcional al tamaño disponible, sino que impone un techo máximo de participación por fuente dentro de cada clase (p. ej., ninguna fuente aporta más del 25-30\% de los ítems de una clase), redistribuyendo la cuota restante entre las demás fuentes de esa clase hasta agotar su disponibilidad. Esto sacrifica algo de eficiencia estadística respecto a un muestreo proporcional puro, a cambio de \textit{prevenir la subrepresentación} de fuentes con menor volumen (La República, Banco de la República). El peso relativo de cada fuente en el corpus recolectado es, al fin y al cabo, un artefacto de cuántas páginas se lograron scrapear de cada una, no una medida de la realidad que valga la pena preservar intacta. El script de muestreo y la tabla de cuotas resultante se registran en \texttt{06\_evaluacion/} junto con el resto del avance del corpus (ver más abajo).

\subsubsection{Etiquetado y validación}

Las noticias verdaderas y falsas reciben su etiqueta de la fuente primaria (medio de referencia o fact-checker IFCN). Las noticias dudosas requieren anotación humano-LLM (no inter-anotador humano clásico, ver limitación metodológica declarada en \texttt{02\_corpus/guia\_anotacion.md}), con medición de acuerdo mediante el coeficiente Kappa de Cohen (meta: $\kappa \geq 0.70$, acuerdo sustancial) entre la etiqueta humana y la del LLM. Los desacuerdos se resuelven con Fable 5.1 como árbitro de razonamiento extendido.

Se implementan controles de calidad ya ejecutados sobre el corpus recolectado hasta ahora, no solo planeados: detección de duplicados mediante similitud de n-gramas (umbral $> 0.85$, con siembra cruzada entre fuentes), un control de calidad manual del 10\% del etiquetado automático (muestra de 25 ítems, 0\% de error encontrado en la primera ejecución), y anonimización pendiente de información personal identificable (campo \texttt{pii\_revisado} en el esquema, ver Marco Normativo).

\section{Proceso metodológico}

\subsection{Fase 1: Construcción del benchmark (Junio--Julio 2026)}

Recolección, etiquetado y validación del corpus. Redacción del marco teórico y estado del arte.

\subsection{Fase 2: Evaluación del LLM como clasificador (Agosto--Septiembre 2026)}

Se evalúan al menos tres LLMs en el benchmark construido, en cuatro configuraciones progresivas:

\begin{enumerate}
    \item \textbf{Nivel 1 -- Zero-shot genérico:} Prompt básico sin contexto de dominio. Constituye el \textit{baseline} ingenuo (equivalente a ``pegar la noticia en ChatGPT'').
    \item \textbf{Nivel 2 -- Zero-shot con prompt de dominio:} Prompt especializado con rol de experto financiero colombiano, instrucciones detalladas sobre indicadores de veracidad y formato de salida estructurado con puntuación de confianza.
    \item \textbf{Nivel 3 -- Few-shot con ejemplos del corpus:} Se proporcionan 3--5 ejemplos reales de cada clase, seleccionados estratégicamente del dominio financiero colombiano.
    \item \textbf{Nivel 4 -- Few-shot con chain-of-thought:} Se solicita al modelo razonar paso a paso: ¿la fuente citada existe?, ¿el dato es verificable?, ¿el lenguaje es neutral o sensacionalista?, ¿menciona entidades regulatorias reales?
\end{enumerate}

Los modelos candidatos incluyen: Gemini Pro/Flash (Google), GPT-4o-mini (OpenAI), y un modelo de código abierto ejecutable en Google Colab (Llama 3.1 8B o Mistral 7B con cuantización).

\subsection{Fase 3: Implementación del agente verificador (Octubre--Noviembre 2026)}

Se implementa el agente basado en ReAct~\cite{Yao2023} con las siguientes herramientas de dominio:

\begin{itemize}
    \item \texttt{consultar\_superfinanciera}: Consulta de comunicados y alertas de la Superintendencia Financiera de Colombia.
    \item \texttt{verificar\_banco\_republica}: Consulta de comunicados oficiales y datos del Banco de la República.
    \item \texttt{buscar\_colombiacheck}: Búsqueda en la base de datos de verificaciones de ColombiaCheck.
    \item \texttt{buscar\_medios\_economicos}: Verificación en La República, Halcones y Palomas y Valora Analitik (Portafolio excluido, ver sección de Adquisición).
    \item \texttt{validar\_fuente\_citada}: Comprobación de la existencia de la fuente citada en la noticia.
\end{itemize}

El agente se integra con el clasificador mediante LangGraph. El flujo de decisión es:

\begin{enumerate}
    \item La noticia ingresa al sistema.
    \item El LLM clasificador emite predicción (verdadera/dudosa/falsa) con puntuación de confianza.
    \item Si la confianza es $\geq 0.80$: se acepta la clasificación del modelo.
    \item Si la confianza es $< 0.80$: se activa el agente verificador.
    \item El agente consulta las fuentes institucionales y emite veredicto con justificación.
    \item El sistema retorna el veredicto final con las fuentes consultadas y el razonamiento.
\end{enumerate}

El umbral no se fija por conveniencia: se determina empíricamente siguiendo el protocolo de clasificación selectiva documentado en \texttt{05\_agente/seleccion\_umbral\_confianza.md} (fundamentado en Chow~\cite{Chow1970} y Geifman y El-Yaniv~\cite{GeifmanElYaniv2017}, ver Marco Tecnológico) -- se barren los valores 0.70, 0.75, 0.80 y 0.85 sobre el split de calibración (\texttt{calibracion\_prompts}, no el de evaluación final, para no sobreajustar el umbral al conjunto que se reporta), se construye la curva riesgo-cobertura, y se elige el umbral según el criterio fijado antes de ver los resultados (maximizar F1 macro del sistema híbrido). Antes de este paso se verifica además que la confianza reportada por el LLM esté razonablemente calibrada (reliability diagram / ECE), ya que no puede asumirse que un LLM reporte una confianza numérica fiable sin verificarlo.

\subsection{Fase 4: Evaluación comparativa y redacción (Diciembre 2026--Enero 2027)}

La evaluación compara cinco configuraciones (los cuatro niveles del clasificador más el sistema híbrido completo) utilizando las siguientes métricas:

\textbf{Métricas primarias:}
\begin{itemize}
    \item F1 macro: promedio de F1 por clase, métrica principal por su robustez ante desbalance.
    \item Precisión y recall por clase: desagregados para verdadera, dudosa y falsa.
    \item Matriz de confusión: visualización de patrones de error.
\end{itemize}

\textbf{Métricas del sistema híbrido:}
\begin{itemize}
    \item Tasa de activación del agente: porcentaje de noticias que activan la verificación.
    \item Mejora marginal: ganancia de precisión cuando el agente interviene.
    \item Latencia: tiempo promedio de procesamiento por noticia.
    \item Cobertura de evidencia: porcentaje de casos donde el agente encontró evidencia relevante.
\end{itemize}

\textbf{Métricas de explicabilidad:}
\begin{itemize}
    \item Fidelidad: consistencia entre la explicación del agente y el veredicto emitido.
    \item Trazabilidad: porcentaje de veredictos con fuentes verificables citadas.
\end{itemize}

\textbf{Significancia estadística:} Se aplica la prueba de McNemar para determinar si las diferencias entre el clasificador solo y el sistema híbrido son estadísticamente significativas ($p < 0.05$).
